\textbf{Racks independientes y elevación RIO-4.}\label{sec:sd4-racks-rio4}
Se seleccionan dos APC NetShelter SX \texttt{AR3104} de 24 U, uno exclusivo de servidores y otro exclusivo de comunicaciones. No se alojan cortafuegos, core o terminaciones WAN en el primero ni servidores en el segundo. La variante publica 600 mm de ancho, 1.070 mm de fondo, 1.198,5 mm de alto, 89,09 kg de masa neta, rieles de 19 pulgadas EIA-310E y profundidad de montaje ajustable de 191 a 914,9 mm. La carga estática publicada es 1.360,8 kg, distinta de la carga dinámica de 1.020,6 kg; ninguna acredita soporte del suelo ni anclaje sísmico. Se amplía comunicaciones de la referencia de 12 U a 24 U sin ampliar el parque de equipos \parencites{p12-apc-ar3104}[RT-06.05, p.~15]{pucv-transversal}.

\begin{table}[htbp]
\centering\small
\caption{Elevación de proyecto RIO-4 por gabinete}\label{tab:sd4-rio4}
\begin{tabularx}{\linewidth}{L{18mm}X X}
\hline
\EncabezadoTabla{Posición} & \EncabezadoTabla{Cómputo, 24 U} & \EncabezadoTabla{Comunicaciones, 24 U} \\ \hline
U1--U2 & PDU A y B, una por U & PDU A y B, una por U \\ \hline
U3--U4 & R360 primero, 1 U; 1 U de holgura & Core primero y segundo, 1 U cada uno \\ \hline
U5--U6 & R360 segundo, 1 U; 1 U de holgura & FG-120G primero y segundo, 1 U cada uno \\ \hline
U7--U8 & Dos U de organización/cableado & Bandejas WAN A y B, 1 U de envolvente cada una \\ \hline
U9--U10 & Libres & Patch panel 1 U y organizador horizontal 1 U \\ \hline
U11--U24 & Libres & Libres \\ \hline
\end{tabularx}
\FuenteTabla{Elevación propia; las holguras y bandejas son reservas de proyecto, no dimensiones atribuidas a equipos sin seleccionar.}
\end{table}

Cómputo reserva 8 U y deja 16 U: $16/24=66,7\,\%$. Comunicaciones reserva 10 U y deja 14 U: $14/24=58,3\,\%$. Las U de holgura se incluyen en ocupación reservada, no se cuentan otra vez como crecimiento. Las U libres reciben paneles ciegos para impedir recirculación, retirados al ampliar; no se añaden ventiladores de gabinete que alteren el balance térmico. El aire entra por el frente y sale por detrás; cables, rieles y PDU conservan el paso de aire, los radios mínimos y la extracción de PSU. El cableado vertical ocupa organizadores laterales y no invade el área útil de montaje. UPS, cadenas de baterías, repuestos y climatización permanecen fuera de estos racks.

La instalación exige montaje fijo sobre apoyos, equipotencialidad, fijación sísmica calculada y sin carga útil sobre ruedas. Se fija una cota propia de 150 kg de equipos/accesorios por rack: masa total de proyecto $89,09+150=239,09$ kg. Sobre $0,600(1,070)=0,642$ m$^2$, la presión media es $239,09(9,80665)/0,642=3,6521$ kN/m$^2$. La obra se calcula para al menos 10 kN/m$^2$ en las zonas TI, además de reacciones concentradas y recorrido de traslado; la placa del gabinete no sustituye esa comprobación. Se admite una profundidad ocupada por equipo de hasta 700 mm, reservando la zona posterior para conectores y cableado sin forzar puertas. La matriz ejecutiva identifica riel/kit, separación entre postes, proyección frontal/posterior, bandeja, masa y radios de cada activo; verifica ReadyRails del R360, ménsulas de core/NGFW y las terminaciones WAN antes de fabricar fijaciones o comprar. La disponibilidad chilena, los kits y la dimensión real de las dos terminaciones no quedan acreditados por la ficha AR3104: si no caben en la elevación y cotas seleccionadas, se modifica la BOM y la elevación manteniendo la separación. Esta dependencia conserva abierto el cierre integral de RIO-4, sin presentar las envolventes como datos de catálogo.

\textbf{Especificación y coordinación de obra OCS-4.}\label{sec:sd4-obra-ocs4}
La obra civil de separación es de cargo del CLIENTE. Synaptix asume su especificación técnica, coordinación y control de interfaces, incluyendo profesionales competentes, sin requerir personal TI propio de la marina. La contraparte del CLIENTE decide y habilita el emplazamiento y contrata/paga la ejecución de obra que le corresponde; esa asignación no exime a Synaptix de definir una solución instalable ni de detectar y corregir incompatibilidades antes de integrar \parencites[RT-06.06, p.~15]{pucv-transversal}[cap.~15, RT-06.01]{caso07}.

El expediente de obra que Synaptix entrega al CLIENTE comprende implantación acotada y accesos; plantas/cortes, separación de zonas, soporte y anclajes; cerramientos BPR-4 y pasos sellados; unifilares A/B, tierra, selectividad y rutas; elevaciones RIO-4 y mapa de enlaces CER-4; ventilación, reposición y drenajes; refrigerante y condiciones de incendio; generación, combustible y contención exterior. Cada documento lleva revisión, responsable profesional, interfaces y lista de materiales/cantidades. Las partidas separan obra del CLIENTE, adquisición de hardware del CLIENTE y servicios de especificación/coordinación/integración de Synaptix. Esa descomposición se concilia en la Oferta Económica sin atribuir un precio o una obra ejecutada a esta memoria.

\begin{table}[htbp]
\centering\small
\caption{Responsabilidades y puntos de liberación de obra OCS-4}\label{tab:sd4-ocs4}
\begin{tabularx}{\linewidth}{L{24mm}X X}
\hline
\EncabezadoTabla{Etapa} & \EncabezadoTabla{Responsabilidad} & \EncabezadoTabla{Condición de avance} \\ \hline
Emplazamiento & CLIENTE habilita acceso/suelo; Synaptix levanta y coordina restricciones con profesionales. & Registro de medidas, servicios, amenazas y emplazamiento aceptado; discrepancias corregidas en plano. \\ \hline
Proyecto ejecutivo & Synaptix especifica y coordina; profesionales competentes suscriben sus especialidades; CLIENTE aprueba la obra a su cargo. & Planos coherentes, permisos/documentación aplicables y cantidades conciliadas; configuración ambiental y eléctrica resuelta antes de liberar compra/obra dependiente. \\ \hline
Ejecución civil & CLIENTE ejecuta mediante contratista competente; Synaptix coordina interfaces y verifica puntos de inspección. & Soporte/anclajes, cerramiento, accesos, rutas y drenajes registrados antes de ocultar o cerrar; cambio autorizado con recálculo previo. \\ \hline
Instalación técnica & Instaladores habilitados ejecutan/certifican sus especialidades; Synaptix integra sin alterar la asignación de coste. & Tierra/protecciones, enlaces, extracción y controles comprobados; planos conforme a obra y expedientes individuales entregados. \\ \hline
Recepción & Synaptix verifica integridad técnica y presenta evidencia; CLIENTE recibe o formula observaciones. & Matriz por recinto/sistema, resultados y pendientes; una no conformidad crítica impide habilitación hasta corregir y repetir la prueba afectada. \\ \hline
\end{tabularx}
\FuenteTabla{Responsabilidades de la oferta y secuencia futura, sin contratos o inspecciones declarados como realizados.}
\end{table}

La recepción civil verifica dimensiones/circulación, accesos independientes, cargas y fijaciones, clasificación del cerramiento/puertas, sellos, rutas separadas y evacuación; la recepción de especialidades agrega electricidad, certificación por enlace, ambiente, combustible e incendio. Un acta civil conforme no acredita autonomía ni climatización. Synaptix conserva registro de cambios, no conformidades, responsables, fechas y pruebas repetidas; entrega planos conforme a obra, manuales y matriz de mantenimiento. La ejecución se coordina con operación de la marina, aislando áreas y servicios afectados, con permisos de trabajo, protección de personas y restitución documentada. Las dimensiones son de instalación nueva; no se atribuyen al edificio existente. OCS-4 mantiene la obligación de especificar/coordinar aunque las condiciones de energía y ambiente aún impidan habilitar.

\textbf{Continuidad individual y ensayo integrado CAI-4.}\label{sec:sd4-autonomia-cai4}
La memoria de 4.2.1 mantiene dos 93T60KMBSB y dos bancos independientes de ochenta HRL12540W por UPS. A plena salida de 60 kW, el cálculo a sesenta minutos es 69,0948 kW AC con factores propios; cada rama se ensaya inicialmente a 60 kW durante sesenta minutos, con generación inhibida, y se mantiene un mínimo de treinta minutos mediante vigilancia y sustitución preventiva. El margen sobre la envolvente de crecimiento inicial de 48,4916 kW, superada por la cota parcial ECI-20 de 50,3600 kW/54,4325 kVA no permite excluir cargas ni compensar una entrada interrumpida. No se suman autonomías A/B ni se transforma la operación sin WAN en horas de batería \parencites{p7-eaton-93t-guia}{csbHRL12540WRA260131}[RT-06.07, p.~15]{pucv-transversal}.

CAI-4 distingue dos fenómenos: al perder la entrada de red/grupo, el inversor online conserva su salida desde el banco; al perder una salida A/B, las cargas deben conservar suministro por su alimentación superviviente. Los R360, FG-120G y core conectan sus dos fuentes a caminos distintos; no se presupone reparto 50/50. Los motores, PC/monitores, CPE, NVR, paneles y convertidores aún deben individualizar entrada, rango admisible y tolerancia temporal. Un ATS/STS de referencia no acredita suministro ininterrumpido: se exige una variante seleccionada cuya transferencia máxima, frente a los modos de falla aplicables y el tiempo de sostenimiento mínimo del activo, no provoque reinicio, pérdida de estado, apertura insegura ni interrupción de captura. Cargas DC deben demostrar fuentes desacopladas o doble entrada y aislamiento; no se conectan dos salidas sin arquitectura autorizada. Se conserva un activo de entrada única como brecha hasta resolver su suministro, aunque la UPS sostenga la energía total.

El mapa de recepción registra por activo y circuito: fuente/PSU, camino A/B y fase, protección/PDU/conversor, modo normal y de rama retirada, carga máxima/transitorio, transferencia o desacoplamiento, continuidad observada y efecto funcional. Prueba la pérdida de entrada AC y de salida por separado, y ambas direcciones de transferencia, con tráfico/captura real, alarma y registro de tiempos. La bancada resistiva de plena carga queda fuera del volumen climatizado, pero los bancos y UPS conservan sus pérdidas interiores reales; en paralelo se ejecuta el ensayo con el parque completo, incluyendo clima, protección y borde, en la banda térmica de autonomía. Se prueba recarga y nuevo corte sin declarar recuperada la reserva hasta medir su condición disponible. La rama en recuperación conserva generación/reabastecimiento y el camino sano conforme a la memoria anterior.

La selección sigue abierta en tres dependencias vinculadas: disipaciones por descarga/recarga y selección del tratamiento ambiental de bancos/UPS; perfil autorizado de recarga, límite por cadena ante pérdida de una de ellas e interfaz autónoma de inhibición por ventilación; BOM de alimentación individual y transferencia de las entradas únicas. Los escenarios de 300/900 m$^3$/h por banco y 9,1917/27,5752 kW latentes agregados no se convierten en equipos de tratamiento ya seleccionados. Añadir ese tratamiento obliga a actualizar carga, fases, grupos y bancos antes de liberar el conjunto. CAI-4 desarrolla la prueba integrada y los límites de selección, pero conserva RT-06.07 por completar; la aceptación futura no sustituye una configuración conjunta todavía indeterminada.
